
\section{$\mathcal M_{\ell+1}(\beta;P_1,\ldots,P_\ell)$ 的定向}
\label{sec:8-4}

\subsection{$\mathcal M_{\ell+1}(\beta;P_1,\ldots,P_\ell)$ 定向的定义}
\label{sec:8-4-1}

设 $L$ 是 $(M,\omega)$ 的相对 spin 拉格朗日子流形。定理
\ref{thm:8-1-1} 说明，从单位圆盘出发、代表同伦类
$\beta\in\pi_2(M,L)$、并在边界上带有 $\ell+1$ 个标记点的伪全纯映射
模空间 $\mathcal M_{\ell+1}^{\mathrm{reg}}(\beta)$ 典范地带有定向。%
\translatornote{源书此处写“定理 8.4.1”，但本小节稍后同号者是定义
8.4.1，并不存在定理 8.4.1；给出上述典范定向的是定理 8.1.1，故改引该
定理。}
讨论定向问题时，固定两个标记点是方便的。固定圆盘上两个点的双全纯
自同构群为 $\mathbb R$，可把它解释为两个全纯圆盘的胶合参数（见第
8.3 节）。一般来说，可以任意选取固定的标记点。例如，若固定
$z_i,z_j$，便可由下列有向空间等式 $(*)$ 定义
$\mathcal M_{\ell+1}^{\mathrm{reg}}(\beta)$ 的定向：
\begin{equation}
 \left\{
 \begin{aligned}
  \widetilde{\mathcal M}_{\ell+1}(\beta)
  &:={}
  \widetilde{\mathcal M}(\beta)
  \times\partial D_0^2\times\cdots\times\partial\check D_i^2
  \times\cdots\times\partial\check D_j^2
  \times\cdots\times\partial D_\ell^2,\\
  \mathcal M_{\ell+1}(\beta)
  &:={}
  \widetilde{\mathcal M}_{\ell+1}(\beta)/\mathbb R.
 \end{aligned}
 \right.
 \tag{*}\label{eq:8-4-star}
\end{equation}
为了定义 $A_\infty$ 运算，只须给
$\mathcal M_{\ell+1}(\beta;P_1,\ldots,P_\ell)$ 的顶维层
$\mathcal M_{\ell+1}^{\mathrm{reg}}(\beta;P_1,\ldots,P_\ell)$ 赋予典范
定向。因此从现在起，我们分别用
$\mathcal M_{\ell+1}(\beta)$ 和
$\mathcal M_{\ell+1}(\beta;P_1,\ldots,P_\ell)$ 简记它们的顶维层。

若以 $\ell+1$ 个字母的对称群 $\mathfrak S_{\ell+1}$ 中的元素
$\sigma$ 改变标记点的次序，它便交换模空间的连通分支。这个微分同胚
保定向，当且仅当 $\sigma$ 是偶置换（见命题~\ref{prop:2-1-22}）。
今后取第 $0$ 个标记点 $z_0$ 与第 $1$ 个标记点 $z_1$ 作为固定点；这与
第 8.3 节的约定相同。若改取其他 $z_i,z_j$，由此得到的
$\mathcal M_{\ell+1}(\beta)$ 的定向，可利用上述
$\mathfrak S_{\ell+1}$-作用同固定 $z_0,z_1$ 的情形作比较。

还应说明只有一个标记点的模空间 $\mathcal M_1(\beta)$。任意添入一个
作为第二个点的标记点，并按上述约定为 $\mathcal M_2(\beta)$ 赋予定向；
再考虑遗忘第二个标记点的映射
\[
 \mathcal M_2(\beta)\longrightarrow\mathcal M_1(\beta).
\]
这个映射的纤维同胚于一个区间，而该区间由 $\partial D^2$ 典范地
定向，因而得到 $\mathcal M_1(\beta)$ 的定向。

现在设 $[P_k,f_k]$ 是 $L$ 中的光滑奇异单形，并且
$\dim P_k=p_k$。置
\[
 \deg P_k=n-p_k,
\]
这正是 $P_k$ 作为上链的次数。由命题~\ref{prop:7-1-1}，
\[
 \operatorname{ev}\colon\mathcal M_{\ell+1}(\beta)
 \longrightarrow L^{\ell+1}
\]
是弱浸没。因此，下面定义~\ref{def:8-4-1} 中的纤维积在约定
\ref{conv:8-2-1}(4) 的意义下是“横截的”。如下定义一个带仓西结构的
空间 $\mathcal M_{\ell+1}(\beta;P_1,\ldots,P_\ell)$。

\begin{definition}{8.4.1}\label{def:8-4-1}
置
\[
\begin{aligned}
 \mathcal M_{\ell+1}(\beta;P_1,\ldots,P_\ell)
 :={}&(-1)^{\epsilon_1}
 \mathcal M_{\ell+1}(\beta)_{
  (\operatorname{ev}_1,\ldots,\operatorname{ev}_\ell)}\\[-2pt]
 &{}\mathop{\times}_{f_1\times\cdots\times f_\ell}
 \left(\prod_{k=1}^{\ell}P_k\right),
\end{aligned}
\]
其中
\[
 \epsilon_1=(n+1)
 \sum_{k=1}^{\ell-1}\sum_{j=1}^{k}\deg P_j.
\]
\end{definition}

\begin{remark}{8.4.2}\label{rem:8-4-2}
利用引理~\ref{lem:8-2-3}(3) 的迭代公式，或注~\ref{rem:8-2-6} 中的
式~\eqref{eq:8-2-7}，可把右端改写成
\[
\begin{aligned}
 &\mathcal M_{\ell+1}(\beta;P_1,\ldots,P_\ell)\\
 &\quad={}
 (-1)^{\sum_{k=1}^{\ell-1}\sum_{j=1}^{k}\deg P_j}\\[-2pt]
 &\qquad{}
 \Bigl(\cdots
   \bigl((\mathcal M_{\ell+1}(\beta)
      {}_{\operatorname{ev}_1}\!\times_{f_1}P_1)
      {}_{\operatorname{ev}_2}\!\times_{f_2}P_2\bigr)
      \mathop{\times}_{f_3}\cdots\Bigr)
      {}_{\operatorname{ev}_\ell}\!\times_{f_\ell}P_\ell.
\end{aligned}
\]
\end{remark}

改变标记点的次序时，定向之差由下面的关键引理给出；它将在后续
论证中起根本作用。

\begin{lemma}{8.4.3}\label{lem:8-4-3}
令 $\sigma=(i,i+1)$ 是 $\ell$ 个字母的对称群
$\mathfrak S_\ell$ 中的对换，其中 $i=1,\ldots,\ell-1$。通过改变标记点
次序，$\sigma$ 在模空间
$\mathcal M_{\ell+1}(\beta;P_1,\ldots,P_i,P_{i+1},\ldots,P_\ell)$
上的作用满足
\[
\begin{aligned}
 &\sigma\bigl(
 \mathcal M_{\ell+1}(\beta;P_1,\ldots,P_i,P_{i+1},\ldots,P_\ell)
 \bigr)\\
 &\qquad={}
 (-1)^{(\deg P_i+1)(\deg P_{i+1}+1)}
 \mathcal M_{\ell+1}
 (\beta;P_1,\ldots,P_{i+1},P_i,\ldots,P_\ell).
\end{aligned}
\]
\end{lemma}

\begin{proof}
由定义，
\begin{equation}
\begin{aligned}
 &\mathcal M_{\ell+1}
  (\beta;P_1,\ldots,P_i,P_{i+1},\ldots,P_\ell)\\
 &\quad={}
 (-1)^{\epsilon_1}
 \mathcal M_{\ell+1}(\beta)_{
 (\operatorname{ev}_1,\ldots,
  \operatorname{ev}_i,\operatorname{ev}_{i+1},\ldots,
  \operatorname{ev}_\ell)}\\[-2pt]
 &\qquad{}\mathop{\times}_{
 f_1\times\cdots\times f_i\times f_{i+1}\times\cdots\times f_\ell}
 \left(\prod_{k=1}^{\ell}P_k\right),
\end{aligned}
\tag{8.4.4.1}\label{eq:8-4-4-1}
\end{equation}
并且
\begin{equation}
\begin{aligned}
 &\mathcal M_{\ell+1}
  (\beta;P_1,\ldots,P_{i+1},P_i,\ldots,P_\ell)\\
 &\quad={}
 (-1)^{\epsilon_2}
 \mathcal M_{\ell+1}(\beta)_{
 (\operatorname{ev}_1,\ldots,
  \operatorname{ev}_i,\operatorname{ev}_{i+1},\ldots,
  \operatorname{ev}_\ell)}\\[-2pt]
 &\qquad{}\mathop{\times}_{
 f_1\times\cdots\times f_{i+1}\times f_i\times\cdots\times f_\ell}
 \left(
  \prod_{k=1}^{i-1}P_k\times P_{i+1}\times P_i
  \times\prod_{k=i+2}^{\ell}P_k
 \right),
\end{aligned}
\tag{8.4.4.2}\label{eq:8-4-4-2}
\end{equation}
其中 $\epsilon_1$ 与定义~\ref{def:8-4-1} 中相同，而
\[
 \epsilon_2
 =\epsilon_1-(n+1)(\deg P_i-\deg P_{i+1}).
\]

为了计算上面两个右端的纤维积定向之差，回忆引理
\ref{lem:8-2-3}(4)。元素 $\sigma$ 通过改变标记点次序作用在
$\mathcal M_{\ell+1}(\beta)$ 上，因而在
$\mathcal M_{\ell+1}(\beta)$ 上诱导一个 $(-1)$-定向同构。显然，
\[
 \prod_{k=1}^{\ell}P_k
 =(-1)^{p_ip_{i+1}}
 \left(
  \prod_{k=1}^{i-1}P_k\times P_{i+1}\times P_i
  \times\prod_{k=i+2}^{\ell}P_k
 \right).
\]
这里 $p_i=\dim P_i$、$p_{i+1}=\dim P_{i+1}$。此外，对纤维积的
底空间定向，有
\[
 \prod_{k=1}^{\ell}L_k
 =(-1)^n
 \left(
  \prod_{k=1}^{i-1}L_k\times L_{i+1}\times L_i
  \times\prod_{k=i+2}^{\ell}L_k
 \right),
\]
其中 $L_k$ 是 $L$ 的一个副本。因此，由引理~\ref{lem:8-2-3}(4)，
\begin{equation}
\begin{aligned}
 &\mathcal M_{\ell+1}(\beta)_{
 (\operatorname{ev}_1,\ldots,
  \operatorname{ev}_i,\operatorname{ev}_{i+1},\ldots,
  \operatorname{ev}_\ell)}
 \mathop{\times}_{
 f_1\times\cdots\times f_i\times f_{i+1}\times\cdots\times f_\ell}
 \left(\prod_{k=1}^{\ell}P_k\right)\\
 &\quad={}
 (-1)^{1+p_ip_{i+1}+n}
 \mathcal M_{\ell+1}(\beta)_{
 (\operatorname{ev}_1,\ldots,
  \operatorname{ev}_i,\operatorname{ev}_{i+1},\ldots,
  \operatorname{ev}_\ell)}\\[-2pt]
 &\qquad{}\mathop{\times}_{
 f_1\times\cdots\times f_{i+1}\times f_i\times\cdots\times f_\ell}
 \left(
  \prod_{k=1}^{i-1}P_k\times P_{i+1}\times P_i
  \times\prod_{k=i+2}^{\ell}P_k
 \right).
\end{aligned}
\tag{8.4.5}\label{eq:8-4-5}
\end{equation}
合并式~\eqref{eq:8-4-4-1}--\eqref{eq:8-4-4-2} 与
\eqref{eq:8-4-5}，可见
\[
 \epsilon_1+\epsilon_2+1+p_ip_{i+1}+n
 \equiv(\deg P_i+1)(\deg P_{i+1}+1)\pmod2,
\]
这就证明了引理~\ref{lem:8-4-3}。
\end{proof}
